% Generated by scripts/gen_blueprint.py from blueprint/nodes/08-roy-lemma.toml. Do not edit.

\chapter{Roy's lemma}
\label{chap:08-roy-lemma}

Roy's lemma on linear spaces of singular matrices. It is the step by which the Structural Rank Conjecture implies the Matrix Coefficient Conjecture of Dasgupta and Kakde \cite{DK24}.

% Prove2Me: Transcendence.singular_matrix_subspace_annihilating_pair -- https://prove2.me/theorems/23adddfc-7432-4ee3-a973-df9f43e8feb0
\begin{lemma}[Roy's lemma]
\label{lem:singular_matrix_subspace_annihilating_pair}
\lean{Transcendence.singular_matrix_subspace_annihilating_pair}
\leanok
Let $F$ be an infinite field and $E$ a linear subspace of the $n\times n$ matrices over $F$, all of whose elements are singular. Then there are non-zero $v,w\in F^{n}$ with
\[ w^{\mathsf T}Av=0\qquad\text{for every }A\in E. \]

{\small\emph{Source:} \cite[Theorem 2.2]{DK24}, with a proof communicated by D.~Roy; a stronger form is \cite[Proposition 12.5]{Wal00}, credited there to Roy (1990).}
\end{lemma}

\begin{proof}
\leanok
Take $A_0\in E$ of maximal rank $r<n$. For every $B\in E$, $B(\ker A_0)\subseteq\operatorname{Im}A_0$: otherwise $A_0+aB\in E$ would have rank $r+1$ for all but finitely many $a\in F$. Take non-zero $v\in\ker A_0$ and non-zero $w$ orthogonal to $\operatorname{Im}A_0$.
\end{proof}
